\section{Four entries under individual proper infiniteness}
\label{sec:four-variables}

Assume in this section that \(C\) satisfies \(\propP\).
We choose four entries that rule out the fixed measurement relations
in Proposition~\ref{prop:graph}.  This will supply the nonvanishing
hypothesis of Lemma~\ref{lem:conditional-shrinking} and hence a
strict contraction between hereditary supports.  A spectral cutoff
is essential: it places a support in an actual operator range, where
the following positivity argument applies.

\begin{lemma}[A spectral cutoff in an actual range]\label{lem:actual-range}
Let \(x,Q\) be bounded positive operators on a Hilbert space
\(\mathcal H\), and put \(U=xQ^2x\).  For every \(\eta>0\),
\begin{equation}\label{eq:actual-range}
 e_{(U-\eta)_+}\mathcal H\subset\ran(xQ),
 \qquad
 \ker Q\cap\ran(xQ)=\{0\}.
\end{equation}
\end{lemma}

\begin{proof}
Let \(f_\eta\) be the bounded Borel function equal to
\(\lambda^{-1}\) for \(\lambda>\eta\) and zero otherwise.
Its norm on the spectrum is at most \(\eta^{-1}\).  If
\(\zeta\in e_{(U-\eta)_+}\mathcal H\), then
\[
 \zeta=Uf_\eta(U)\zeta
      =xQ\bigl(Qx f_\eta(U)\zeta\bigr),
\]
which proves the actual-range inclusion.  If
\(\zeta=xQ\xi\) and \(Q\zeta=0\), then
\[
 0=\langle QxQ\xi,\xi\rangle=\|x^{1/2}Q\xi\|^2,
\]
so \(\zeta=xQ\xi=0\).  The assertion concerns the actual range;
no closed-range conclusion is being used.
\end{proof}

\begin{proposition}[Four entries avoid the fixed measurements]
\label{prop:four-avoidance}
Suppose \(C\) satisfies \(\propP\).  Let \(X\in C_+\) and
\(B_1,B_2,B_3,B_4\in C\).  With
\[
 \alpha_0=\frac1{32}<\alpha=\frac1{16}
 <\beta=\frac18<s=\frac14,\qquad h=\frac12,
\]
there exist contractions \(T_1,\ldots,T_4\) in
\(\Her((X-\alpha_0)_+)\) for which no representation \(\pi\)
of \(C\) satisfies
\begin{equation}\label{eq:four-forbidden}
 \|\pi(X)\|>s,\qquad \pi(B_i)=\pi(T_i)\quad(1\le i\le4).
\end{equation}
\end{proposition}

\begin{proof}
We argue by contradiction.  If every four-entry choice could be
matched in a representation, the fixed elements could be passed to
a norm ultraproduct.  In an irreducible representation of the
resulting image of \(C\), the first pair would give an isometry with
nonzero defect, while the last pair would rule out that defect.

All formulas described as fixed below depend only on \(X,B_i\)
and the displayed constants.  Define
\begin{equation}\label{eq:four-first-fixed}
 x=(X-\alpha)_+,\qquad V=(X-\beta)_+,\qquad
 Q=(h1-B_1B_1^*)_+,\qquad U=xQ^2x.
\end{equation}
Here \(Q\in\widetilde C_+\), whereas \(x,V,U\in C_+\).
Let \(\fin(C)\) denote the set of finite subsets of \(C\), and order
\[
 \Lambda=\mathbb N_{\ge1}\times\fin(C),
\]
by the usual order in the first coordinate and inclusion in the
second.  At \(\lambda=(m,E)\), put \(\eta=1/m\), so
\(\eta\to0\) along this directed set, and every fixed \(w\in C\)
eventually belongs to \(E\).

\medskip\noindent\emph{The first two entries.}
The first pair will force \(UwV=0\) for every fixed \(w\in C\)
in a limiting image of \(C\).  At the index \((m,E)\), we test
only the finitely many \(w\in E\).  A positive sum lets one norm
bound control all of these tests: set
\begin{equation}\label{eq:four-D}
 A_\eta=(U-\eta)_+,\qquad
 D_0=A_\eta\left(\sum_{w\in E}wVw^*\right)A_\eta,\qquad
 D=g_\eta(D_0),
\end{equation}
where \(g_\eta\) is a ramp that is zero on \([0,\eta]\) and
one on \([2\eta,\infty)\).  In fact,
\[
 D_0=\sum_{w\in E}
 (A_\eta wV^{1/2})(A_\eta wV^{1/2})^*.
\]
Thus a small norm for \(D_0\) makes each
\(A_\eta wV^{1/2}\) small.  We have
\[
 D_0\in\Her(A_\eta)\cap\Id(V),\qquad
 e_D\le e_{(D_0-\eta)_+}\le e_{A_\eta}\le e_x.
\]
Since \(V\) is properly infinite, ideal comparison gives
\(D_0\oplus D_0\precsim V\).  Apply
Lemma~\ref{lem:transport} with cutoff \(\eta\), split the
carrier into two columns, and multiply these columns by
\(D^{1/2}\).  Taking adjoints of the resulting two elements gives
contractions \(T_1,T_2\in C\) with
\begin{equation}\label{eq:four-first-entries}
 T_iT_j^*=\delta_{ij}D,\qquad e_VT_i^*=T_i^*
 \quad(i,j=1,2).
\end{equation}
Their left supports lie under \(e_D\) and their right supports
under \(e_V\).  Both are under \(e_x\le e_{(X-\alpha_0)_+}\);
they therefore belong to the required hereditary algebra.

Consider a representation satisfying
\eqref{eq:four-forbidden} for the entries at this index, and
omit its symbols.  Then \(B_1B_1^*=D\) and
\(Q=(h1-D)_+\).  The spectral subspace of \(D\) for
\((h,\infty)\) is killed by \(Q\) and lies in
\(e_D\mathcal H\subset e_{(U-\eta)_+}\mathcal H\).
Lemma~\ref{lem:actual-range} shows that this subspace is zero.
Thus
\begin{equation}\label{eq:four-first-bounds}
 B_1B_1^*=D\le h1,\qquad \|D_0\|\le2\eta
 \quad\text{in this representation}.
\end{equation}
For the second conclusion, if the norm of \(D_0\) were at
least \(2\eta\), spectral mapping and the plateau of \(g_\eta\)
would give \(\|D\|=1>h\).

\medskip\noindent\emph{A fixed prospective defect and the last two entries.}
Choose ramps
\[
 H=h_0(X),\qquad K=k_0(X)
\]
so that \(h_0\) vanishes on a neighborhood of
\([0,\alpha]\) and equals one on \([\beta,\infty)\), and
\(k_0\) vanishes on a neighborhood of \([0,\alpha_0]\) and
equals one on \([\alpha,\infty)\).  In particular \(KH=H\).
Set
\begin{equation}\label{eq:four-r}
 r_i=h^{-1/2}B_i^*H\quad(i=1,2),\qquad
 b_0=(1-H^2)^{1/2}.
\end{equation}
The reason for these definitions is the following conditional
calculation.  In the limiting irreducible representation below, the
first pair will give
\[
 r_i^*r_j=\delta_{ij}H^2,\qquad Hr_i=r_i\quad(i,j=1,2),
 \qquad H\ne0.
\]
These relations are not being asserted in \(C\) at this stage.
They imply \(b_0r_i=0\), and hence
\[
 v=r_1+b_0,\qquad v^*v=H^2+b_0^2=1,\qquad v^*r_2=0.
\]
Since \(r_2^*r_2=H^2\ne0\), the isometry \(v\) has a nonzero
defect projection \(1-vv^*\).  This is the scaling-element
construction of Blackadar--Cuntz \cite{BlackadarCuntz1982}, in the
form recalled in \cite[Remark~2.4]{PasnicuRordam2007}; the second
entry supplies the additional nonzero orthogonal range.

To choose the last two entries before taking any limit, we encode
this prospective defect by one fixed element of \(C\).  Compression
by \(K\) puts the formula in the required hereditary algebra, and
clipping makes it a positive contraction even before the preceding
relations hold.  Let
\(\chi(t)=\min\{1,\max\{0,t\}\}\) on \(\mathbb R\), and define
\begin{equation}\label{eq:four-P0}
 P_0=\chi\!\left(
 K\bigl[1-(r_1+b_0)(r_1+b_0)^*\bigr]K
 \right)\in C.
\end{equation}
The compressed selfadjoint element is in \(C\) and supported
under \(e_K\); since \(\chi(0)=0\), continuous functional
calculus in \(\Her(K)\) gives
\[
 e_{P_0}\le e_K\le e_{(X-\alpha_0)_+}.
\]
This membership and support statement uses no relation from a
representation.

Define the following additional fixed elements:
\begin{equation}\label{eq:four-second-fixed}
 x'=(P_0-1/2)_+,\qquad
 Q'=(h1-B_3B_3^*)_+,\qquad R=x'(Q')^2x'.
\end{equation}
At the same index \(\eta\), property \(\propP\) gives
\[
 ((R-\eta)_+)^{\oplus2}\precsim(R-\eta)_+.
\]
Apply Lemma~\ref{lem:transport} to this comparison with cutoff
\(\eta\).  If \(t_R\) is a ramp of \(R\),
zero on \([0,2\eta]\) and one on \([4\eta,\infty)\), the two
carrier columns times \(t_R\) give contractions \(T_3,T_4\in C\)
such that
\begin{equation}\label{eq:four-last-entries}
 T_3^*T_4=0,\qquad T_3^*T_3=T_4^*T_4=t_R^2.
\end{equation}
Their left supports lie under \(e_{(R-\eta)_+}\) and their
right supports under \(e_{(R-2\eta)_+}\).  Both are under
\(e_{x'}\le e_{(X-\alpha_0)_+}\).  Transport preserves the
represented norms, so
\begin{equation}\label{eq:four-last-norm}
 \|\pi(R)\|>4\eta\quad\Longrightarrow\quad
 \|\pi(T_3)\|=1.
\end{equation}

In a representation satisfying \eqref{eq:four-forbidden},
\(B_3=T_3\).  Apply Lemma~\ref{lem:actual-range} to \(x',Q',R\).
The spectral subspace of \(B_3B_3^*\) above \(h\) is killed by
\(Q'\) and lies under \(e_{(R-\eta)_+}\), so it is zero.
Equation~\eqref{eq:four-last-norm} then gives
\begin{equation}\label{eq:four-last-bounds}
 B_3B_3^*\le h1,\qquad \|R\|\le4\eta
 \quad\text{in this representation}.
\end{equation}

\medskip\noindent\emph{Passing only fixed elements to an ultraproduct.}
Suppose, for a contradiction, that every index in \(\Lambda\)
admits a representation satisfying \eqref{eq:four-forbidden}.
Choose one, denoted \(\pi_\lambda\), on its nondegenerate
essential subspace.  There is no size issue in this choice: from
any such representation, choose a vector detecting
\(\|\pi_\lambda(X)\|>s\) and restrict to its reducing cyclic
subrepresentation.  A set of cyclic representations suffices.
Choose an ultrafilter \(\omega\) containing all tails of
\(\Lambda\), and form the norm ultraproduct
\[
 \mathcal B=\prod_{\omega}B(\mathcal H_\lambda)
 =\frac{\prod_\lambda B(\mathcal H_\lambda)}
 {\{(z_\lambda):\lim_{\omega}\|z_\lambda\|=0\}}.
\]
The norm of a class is \(\lim_{\omega}\|z_\lambda\|\).
Let \(\Phi:C\to\mathcal B\) be
\(\Phi(z)=[(\pi_\lambda(z))]\), and write \(B=\Phi(C)\).
The image is a closed \(C^*\)-subalgebra.  The unital extensions
carry the unitized formulas into \(C^*(B,1_{\mathcal B})\).
We suppress \(\Phi\) on fixed elements from now on.

For any fixed \(w\in C\), eventually \(w\in E\).  Positivity
of the summands in \eqref{eq:four-D} and
\eqref{eq:four-first-bounds} yield
\[
 \|\pi_\lambda((U-\eta)_+wV^{1/2})\|^2\le2\eta
\]
on that tail.  Since \(\|U-(U-\eta)_+\|\le\eta\),
\begin{equation}\label{eq:four-UCV-estimate}
 \|\pi_\lambda(UwV)\|
 \le\sqrt{2\eta}\,\|V\|^{1/2}
      +\eta\|w\|\|V\|\longrightarrow0.
\end{equation}
It follows that \(UBV=0\) in the image algebra.  The other exact
relations and bounds at each index give
\begin{equation}\label{eq:four-ultraproduct}
 \begin{gathered}
 \|X\|\ge s,\qquad B_1B_1^*\le h1,\qquad
 B_iB_j^*=\delta_{ij}B_1B_1^*,\quad HB_i^*=B_i^*
       \quad(i,j=1,2),\\
 UBV=0,\qquad R=0,\qquad B_3B_3^*\le h1,\qquad
 B_3^*B_4=0,\quad B_3^*B_3=B_4^*B_4.
 \end{gathered}
\end{equation}
For \(HB_i^*=B_i^*\), use \(e_VT_i^*=T_i^*\) and
\(H=1\) on the support of \(V\).
Fix also a continuous \(h_P:[0,1]\to[0,1]\) with
\(h_P(0)=0\) and \(h_P=1\) on \([1/2,1]\).  The left
support of \(T_4\) lies under \(e_{x'}\), so we obtain
\begin{equation}\label{eq:four-B4-support}
 h_P(P_0)B_4=B_4.
\end{equation}
Every element in \eqref{eq:four-ultraproduct} and
\eqref{eq:four-B4-support} is the image of a fixed algebra
element or of its fixed unitized functional calculus.  We have
not placed the varying \(T_i,D,D_0,t_R\) in \(B\), and we do
not require proper infiniteness of \(B\) or of \(\mathcal B\).
The directed finite sets in \eqref{eq:four-UCV-estimate} handle
all \(w\in C\) without a separability assumption.

\medskip\noindent\emph{A nonzero defect gives the contradiction.}
Since \(\|X\|\ge s>\beta\), the image \(V=(X-\beta)_+\)
is nonzero.  Choose an irreducible representation of \(B\)
in which \(V\ne0\).  Extend this nondegenerate representation
to \(C^*(B,1_{\mathcal B})=B+\mathbb C1_{\mathcal B}\) by
sending \(1_{\mathcal B}\) to the identity operator; this remains
valid if \(B\) has a unit different from \(1_{\mathcal B}\).
Thus the unitized expressions \(Q,b_0,Q'\) and the defect below
have their stated represented meanings.  Again suppress the
representation symbols.  The closed span of \(BV\mathcal H\) is a nonzero
reducing subspace, and hence all of \(\mathcal H\).
The relation \(UBV=0\) therefore gives \(U=0\).
Put \(D_1=B_1B_1^*\).  Since \(D_1\le h1\), we have
\(Q=h1-D_1\), and
\[
 0=U=(Qx)^*(Qx)
 \quad\Longrightarrow\quad D_1x=hx.
\]
By closure \(D_1=h1\) on \(e_x\mathcal H\).  The support of
\(H\) lies there.  Equations~\eqref{eq:four-r} and
\eqref{eq:four-ultraproduct} give
\begin{equation}\label{eq:four-isometry-relations}
 r_i^*r_j=\delta_{ij}H^2,\qquad Hr_i=r_i
 \quad(i,j=1,2).
\end{equation}
Also \(HV=V\ne0\), so \(H\ne0\) and \(r_2\ne0\).

The conditional calculation preceding \eqref{eq:four-P0} now
applies: \(v=r_1+b_0\) is an isometry with \(v^*r_2=0\).
Its defect projection
\[
 P=1-vv^*
\]
is nonzero, since \(Pr_2=r_2\ne0\).

We check that this projection is the represented fixed element
\(P_0\), a point for which both sides of the cutoff matter.
The scalar ramps give \(KH=H\).  From \(Hr_1=r_1\) follows
\(Kr_1=r_1\), and the definition \(r_1=h^{-1/2}B_1^*H\)
gives \(r_1K=r_1\).  Also \(K\) commutes with \(b_0\).
Expanding the defect gives
\[
 P=H^2-r_1r_1^*-r_1b_0-b_0r_1^*.
\]
Each term is fixed by multiplication by \(K\) on both sides.
Therefore \(KPK=P\), and \(\chi(P)=P\) because \(P\) is a
projection.  Formula~\eqref{eq:four-P0} yields
\begin{equation}\label{eq:four-P0-is-P}
 P_0=P.
\end{equation}

Now \(x'=(P_0-1/2)_+=P/2\).  Put \(D_3=B_3B_3^*\le h1\);
then \(Q'=h1-D_3\).  The relation \(R=0\) gives
\(Q'x'=0\), so
\begin{equation}\label{eq:four-last-contradiction}
 D_3P=hP.
\end{equation}
Equations~\eqref{eq:four-B4-support} and
\eqref{eq:four-P0-is-P} say \(B_4=PB_4\).
But \(B_3^*B_4=0\), so \(D_3B_4=0\), whereas
\eqref{eq:four-last-contradiction} gives \(D_3B_4=hB_4\).
Thus \(B_4=0\), and
\(B_3^*B_3=B_4^*B_4\) gives \(B_3=0\).
This contradicts \eqref{eq:four-last-contradiction}, since
\(h>0\) and \(P\ne0\).

The contradiction shows that some index has no representation
satisfying \eqref{eq:four-forbidden}.  The four entries chosen
at that index prove the proposition.
\end{proof}

We now apply the avoidance proposition to the graph measurements.
This gives the form of shrinking that will be used throughout the
next section.

\begin{proposition}[Shrinking under individual proper infiniteness]
\label{prop:shrinking-P}
Suppose \(C\) satisfies \(\propP\).  Let \(a,y\in C\) be positive
contractions and let \(c\in\widetilde C\) be a contraction.  There
are positive contractions \(p_1\in\Her(a)\) and
\(q_1\in\Her(y)\), with \(p_1\) rich, such that for every
representation \(\pi\) of \(C\),
\begin{equation}\label{eq:shrinking-P-signals}
 \|\pi(a)\|>1/2\ \Longrightarrow\ \|\pi(p_1)\|=1,
 \qquad
 \|\pi(y)\|>1/2\ \Longrightarrow\ \|\pi(q_1)\|=1,
\end{equation}
and, in the universal bidual,
\begin{equation}\label{eq:shrinking-P-support}
 \|e_{p_1}c e_{q_1}\|<1.
\end{equation}
\end{proposition}

\begin{proof}
Every positive contraction of \(C\) is rich under \(\propP\), since
each of its scalar cutoffs is properly infinite.  In
Proposition~\ref{prop:graph}, fix \(k=4\) before the measurements
are computed, and choose, for example,
\(\gamma_g=1/2\) and \(\sigma_g=1/(2\sqrt5)\).  Then
\(\gamma_g^2+4\sigma_g^2<1\).  The graph proposition gives fixed
\(X\) and \(B_1,\ldots,B_4\) for the inputs \(a,y,c\).
Apply Proposition~\ref{prop:four-avoidance} to these elements and
use its contractions \(T_i\) as the graph choices \(S_i\).
The resulting rich positive contraction \(p\in\Her(a)\) has the
required norm-one signal for \(a\).

Put \(d=y(1-c^*p^2c)y\).  If a representation \(\pi\) had
\(\pi(d)=0\) and \(\|\pi(y)\|>1/4\), the graph implication would
give \(\|\pi(X)\|>1/4\) and \(\pi(B_i)=\pi(T_i)\) for every
\(i\).  These are exactly the relations excluded by
Proposition~\ref{prop:four-avoidance}.  Thus
\begin{equation}\label{eq:nonvanishing-P}
 \pi\bigl(y(1-c^*p^2c)y\bigr)=0
 \quad\Longrightarrow\quad \|\pi(y)\|\le1/4
\end{equation}
for every representation.  The same \(p\) is rich and carries the
\(a\)-signal, so Lemma~\ref{lem:conditional-shrinking} applies to
\(a,y,c,p\) and proves the stated conclusions.  This argument also
covers zero inputs.
\end{proof}
